% theme: proofs_with_inscribed_angles
% 表：代表的な穴埋め証明4題。裏：完全記述3題と、その証明を利用する計算。
\MajorQuestionLesson{円周角の定理の証明問題}
\SetRowSpace{3mm}

\TypeQuestion{図で、$\Ang{ACB}=\Ang{ADB}$ である。4点 $\mathrm{A},\mathrm{B},\mathrm{C},\mathrm{D}$ が同一円周上にあることを示す証明の空欄をうめなさい。}{converse_concyclic_fill}
\FigProblemBlock{39mm}{fig/lesson_04_concyclic_01}{
\WholeQuestion{点 $\mathrm{A},\mathrm{B},\mathrm{C}$ を通る円をかく。この円において、弦 $\mathrm{AB}$ に対する円周角は $\Ang{ACB}$ である。仮定より、$\Ang{ADB}=\Blank{ア}$ であるから、円周角の定理の逆より、点 $\mathrm{D}$ はこの円の\Blank{イ}にある。したがって、4点 $\mathrm{A},\mathrm{B},\mathrm{C},\mathrm{D}$ は同一円周上にある。}{}}
{} 

\TypeQuestion{図で、$\Ang{ABC}+\Ang{ADC}=180^\circ$ である。4点 $\mathrm{A},\mathrm{B},\mathrm{C},\mathrm{D}$ が同一円周上にあることを示す証明の空欄をうめなさい。}{opposite_angles_concyclic_fill}
\FigProblemBlock{39mm}{fig/lesson_04_concyclic_02}{
\WholeQuestion{点 $\mathrm{A},\mathrm{B},\mathrm{C}$ を通る円と、直線 $\mathrm{CD}$ とのもう1つの交点を $\mathrm{D}'$ とする。円に内接する四角形の向かい合う角の和は $180^\circ$ だから、$\Ang{ABC}+\Ang{AD'C}=180^\circ$ である。仮定と比べると、$\Ang{ADC}=\Blank{ア}$ となるので、点 $\mathrm{D}$ と点 $\mathrm{D}'$ は一致する。よって、\Blank{イ}。}{}}
{} 

\TypeQuestion{図で、2本の弦 $\mathrm{AC},\mathrm{BD}$ は点 $\mathrm{P}$ で交わっている。$\Tri{APD}\Similar\Tri{BPC}$ を示す証明の空欄をうめなさい。}{crossing_chords_similarity_fill}
\FigProblemBlock{39mm}{fig/lesson_04_similarity_01}{
\WholeQuestion{$\Tri{APD}$ と $\Tri{BPC}$ において、同じ弧 $\mathrm{CD}$ に対する円周角は等しいから、$\Ang{CAD}=\Ang{CBD}$。したがって、$\Ang{PAD}=\Blank{ア}$。また、対頂角は等しいから、$\Ang{APD}=\Blank{イ}$。よって、\Blank{ウ}ので、$\Tri{APD}\Similar\Tri{BPC}$。}{}}
{} 

\TypeQuestion{四角形 $\mathrm{ABCD}$ は円に内接し、辺 $\mathrm{AD},\mathrm{BC}$ の延長の交点を $\mathrm{E}$ とする。$\Tri{ABE}\Similar\Tri{CDE}$ を示す証明の空欄をうめなさい。}{cyclic_exterior_similarity_fill}
\FigProblemBlock{39mm}{fig/lesson_04_similarity_02}{
\WholeQuestion{$\Tri{ABE}$ と $\Tri{CDE}$ において、円に内接する四角形の外角は、その向かいの内角に等しいから、$\Ang{BAE}=\Blank{ア}$。また、$\Ang{AEB}=\Blank{イ}$。よって、\Blank{ウ}ので、$\Tri{ABE}\Similar\Tri{CDE}$。}{}}
{} 

\AnswerColumnBreak
\LessonQuestionBreak

\TypeQuestion{図で、$\Ang{ACB}=\Ang{ADB}$ である。}{converse_concyclic_proof}
\FigProblemBlock{43mm}{fig/lesson_04_concyclic_03}{
\QQ{4点 $\mathrm{A},\mathrm{B},\mathrm{C},\mathrm{D}$ が同一円周上にあることを証明しなさい。}{}
\QQ{$\Ang{BAC}=48^\circ$、$\Ang{BCD}=71^\circ$ のとき、$\Ang{CAD}$ の大きさを求めなさい。}{}
}{}

\TypeQuestion{図で、2本の弦 $\mathrm{AC},\mathrm{BD}$ は点 $\mathrm{P}$ で交わっている。}{crossing_chords_similarity_proof}
\FigProblemBlock{43mm}{fig/lesson_04_similarity_03}{
\QQ{$\Tri{APD}\Similar\Tri{BPC}$ となることを証明しなさい。}{}
\QQ{$\Lseg{AP}=6\mathrm{cm}$、$\Lseg{BP}=4\mathrm{cm}$、$\Lseg{CP}=9\mathrm{cm}$ のとき、$\Lseg{DP}$ の長さを求めなさい。}{}
}{}

\TypeQuestion{四角形 $\mathrm{ABCD}$ は円に内接し、辺 $\mathrm{AD},\mathrm{BC}$ の延長の交点を $\mathrm{E}$ とする。}{cyclic_exterior_similarity_proof}
\FigProblemBlock{43mm}{fig/lesson_04_similarity_04}{
\QQ{$\Tri{ABE}\Similar\Tri{CDE}$ となることを証明しなさい。}{}
\QQ{$\Lseg{AB}=8\mathrm{cm}$、$\Lseg{CD}=6\mathrm{cm}$、$\Lseg{DE}=9\mathrm{cm}$ のとき、$\Lseg{BE}$ の長さを求めなさい。}{}
}{}
